\section{Simulation Setup}
\label{app:simulation-setup}

Section~\ref{sec:experimental-setup} defines the common filter, sensor, communication, operating-mode, and Monte Carlo settings. The remaining constants needed to reproduce the scenario are the formation geometry, target initialization, and paired-run handling below.

The eight sensors form two four-sensor groups. Their formation centers start at $[-80,35]^\mathsf{T}$ and $[-80,-35]^\mathsf{T}$ and move with nominal velocity $[0.8,0]^\mathsf{T}$. Each group follows a four-agent leader-follower pattern with spacing $20$. The communication graph is fully connected inside each group, with cross links $1\leftrightarrow5$, $2\leftrightarrow6$, $3\leftrightarrow7$, and $4\leftrightarrow8$; the communication range is $150$.

The ten targets are initialized as the three staggered groups summarized in Section~\ref{sec:experimental-setup}. Their group centers are $[70,80]^\mathsf{T}$, $[80,0]^\mathsf{T}$, and $[70,-80]^\mathsf{T}$, with spacings $30$, $25$, and $20$, respectively. Births start at time step one, repeat every eight steps, and last through the 100-step scenario.

Within each deterministic seed, all compared arms share the same generated truth, measurements, communication realization, and delivered measurement sets. The ideal-communication comparison reuses the same geometry and targets but sets communication level $0$, uses effectively unlimited delivery, and sets $p_{\mathrm{drop}}=0$. The appendix-only AA diagnostics use the same communication-constrained scenario, with a shorter AA horizon and stronger AA-specific pruning to control mixture growth.
